\documentclass{article}
\usepackage[utf8]{inputenc}
\usepackage{amsmath}
\usepackage{geometry}
\geometry{margin=2cm}
\begin{document}

\section*{Resolução da Questão 94 -- ENEM 2024 -- Ciências da Natureza}

\subsection*{Interpretação do enunciado, estratégia e análise}
A questão apresenta a hidroxiapatita, Ca$_5$(PO$_4$)$_3$OH, como componente do esmalte dentário e um equilíbrio químico no meio bucal envolvendo reações de mineralização e desmineralização. O objetivo é identificar, com base no efeito do pH dos cremes dentais comerciais, qual deles promove maior desmineralização do esmalte.

A estratégia para resolução é analisar como o pH, indicador da concentração de íons H$^{+}$, afeta o equilíbrio químico da hidroxiapatita em solução aquosa. Sabemos que aumento de H$^{+}$ (meio ácido) desloca o equilíbrio para a dissolução do mineral, favorecendo a desmineralização.

\subsubsection*{Análise técnica}
A reação de equilíbrio pode ser representada por:
\[
\text{Ca}_5(\text{PO}_4)_3\text{OH (s)} + \text{H}^+ (\text{aq}) \rightleftharpoons 5\,\text{Ca}^{2+} (\text{aq}) + 3\,\text{PO}_4^{3-} (\text{aq}) + \text{H}_2\text{O} (\text{l})
\]

Um meio com maior concentração de H$^{+}$ (pH menor) favorece a reação à direita (dissolução), agravando a desmineralização e consequentemente a fragilidade do esmalte dentário.

A tabela abaixo apresenta os pHs dos cremes dentais:

\begin{center}
\begin{tabular}{|c|c|c|}
\hline
\textbf{Creme dental} & \textbf{Agente de polimento} & \textbf{pH} \\
\hline
I & Bicarbonato de sódio & 9,5 \\
II & Carbonato de cálcio & 11,0 \\
III & Citrato de potássio & 7,7 \\
IV & Dióxido de silício & 6,9 \\
V & Fosfato de cálcio & 7,3 \\
\hline
\end{tabular}
\end{center}

Nota-se que o creme IV apresenta o menor pH, indicando o ambiente mais ácido entre as opções e, assim, o maior potencial para causar desmineralização da hidroxiapatita.

\section*{Conteúdo teórico e mini revisão}
\subsection*{Equilíbrio químico}
O equilíbrio químico ocorre quando as velocidades da reação direta e inversa se igualam, mantendo constantes as concentrações dos reagentes e produtos.

\subsection*{pH e concentração de íons H$^{+}$}
O pH é definido por:
\[
pH = -\log[\text{H}^{+}]
\]
Um pH menor que 7 indica meio ácido, ou seja, maior concentração de íons H$^{+}$, que pode afetar o equilíbrio químico deslocando-o no sentido que consome os íons H$^{+}$.

\subsection*{Desmineralização do esmalte dentário}
A hidroxiapatita, principal componente do esmalte, se dissolve com a presença elevada de íons H$^{+}$, causando a fragilização dos dentes e facilitando o aparecimento de cáries.

\section*{Resolução detalhada e "pulo do gato"}
Com base na reação:
\[
\text{Ca}_5(\text{PO}_4)_3\text{OH (s)} + \text{H}^+ (\text{aq}) \rightleftharpoons 5\,\text{Ca}^{2+} (\text{aq}) + 3\,\text{PO}_4^{3-} (\text{aq}) + \text{H}_2\text{O} (\text{l})
\]

A presença de mais H$^{+}$ desloca o equilíbrio para a direita, aumentando a desmineralização.

Analisando os pHs, o creme dental IV, com pH 6,9, é o mais ácido, portanto, será o que mais favorece essa desmineralização.

\textbf{Conclusão:} A alternativa correta é a letra \textbf{D}.

\textit{Pulo do gato:} Para identificar o creme que promove maior desmineralização, basta encontrar o de menor pH na tabela, pois a baixa acidez é que protege o esmalte, enquanto o meio ácido promove dissolução.

\section*{Armadilhas comuns e dicas para o ensino}
\subsection*{Armadilhas comuns}
\begin{itemize}
  \item Assumir que pH maior (meio alcalino) provoca mais desmineralização, quando na verdade protege o esmalte.
  \item Confundir o agente de polimento (parte física) com o efeito químico do pH no equilíbrio.
  \item Subestimar o efeito de um pH levemente ácido (como 6,9) sobre a dissolução do esmalte.
  \item Generalizar que qualquer mudança no pH promove desmineralização, sem analisar o sentido do deslocamento do equilíbrio.
\end{itemize}

\subsection*{Dicas para ensinar}
\begin{itemize}
  \item Explicitar o conceito de equilíbrio químico e como o excesso de H$^{+}$ influencia a dissolução dos minerais.
  \item Reforçar a relação entre pH e concentração de íons H$^{+}$ e explicar a classificação de ácidos e bases.
  \item Utilizar tabelas comparativas para mostrar os efeitos práticos do pH em diferentes contextos.
  \item Mostrar o impacto da química para a saúde, contextualizando assim maior interesse.
  \item Estimular o raciocínio crítico e lógico para interpretar efeitos químicos em vez de decorar valores.
\end{itemize}

\section*{Código da questão}
CN\_DENTAL\_EQ\_001

\section*{Anexos: Imagens do enunciado}
\begin{center}
\textbf{Reação química representando o equilíbrio da hidroxiapatita}

\fbox{\parbox{10cm}{
\[
\text{Ca}_5(\text{PO}_4)_3\text{OH (s)} + \text{H}^+ (\text{aq}) \rightleftharpoons 5\,\text{Ca}^{2+} (\text{aq}) + 3\,\text{PO}_4^{3-} (\text{aq}) + \text{H}_2\text{O} (\text{l})
\]
}}

\vspace{1cm}

\textbf{Tabela de cremes dentais, agentes e pH}

\fbox{\parbox{14cm}{
\begin{tabular}{|c|c|c|}
\hline
\textbf{Creme dental} & \textbf{Agente de polimento} & \textbf{pH} \\
\hline
I    & Bicarbonato de sódio & 9,5  \\
II   & Carbonato de cálcio    & 11,0 \\
III  & Citrato de potássio   & 7,7  \\
IV   & Dióxido de silício    & 6,9  \\
V    & Fosfato de cálcio     & 7,3  \\
\hline
\end{tabular}
}}
\end{center}


\end{document}